% !PRINTED 293
\weberpage{293}

In this case, since the signs of $\alpha$, $\beta$,
$\gamma$, $\delta$ may all be reversed at the same time, we may assume:
\[
\alpha+\delta\equiv 2, \quad \alpha\delta-\beta\gamma\equiv 1 \quad (\mathrm{mod}\ p),
\]
and thus we obtain from (8):
\[
\begin{aligned}
\alpha_2 &\equiv -1+2\alpha, & \beta_2 &\equiv 2\beta\\
\delta_2 &\equiv -1+2\delta, & \gamma_2 &\equiv 2\gamma\\
\alpha_3 &\equiv -2+3\alpha, & \beta_3 &\equiv 3\beta\\
\delta_3 &\equiv -2+3\delta, & \gamma_3 &\equiv 3\gamma
\end{aligned}\ \ (\mathrm{mod}\ p),
\]
from which, by the inference from $m$ to $m+1$, it follows that:
\weq{9}{\begin{aligned}
\alpha_m &\equiv -(m-1)+m\alpha, & \beta_m &\equiv m\beta\\
\delta_m &\equiv -(m-1)+m\delta, & \gamma_m &\equiv m\gamma
\end{aligned}\ \ (\mathrm{mod}\ p).}
\gesperrt{From this one sees that all the $A^m$ again
belong to Case I., and that $p$ is the degree of $A$.}

The analogous consideration of the two other cases has been carried out by
\gesperrt{Serret}, but appears dispensable for our purpose.

On the other hand, we wish to add here the further theorem, \gesperrt{that
the whole group $\mathfrak{L}_0$ can be composed out of the
two following special substitutions}
\weq{10}{A=\begin{pmatrix}1, & 0\\ 1, & 1\end{pmatrix}, \quad B=\begin{pmatrix}0, & 1\\ -1, & 0\end{pmatrix}.}
The proof follows from \S\,30, if one observes that
for every system of numbers $\alpha$, $\beta$, $\gamma$, $\delta$ satisfying the condition $\alpha\delta-\beta\gamma\equiv 1\ (\mathrm{mod}\ p)$
the system of numbers $\alpha'$, $\beta'$, $\gamma'$, $\delta'$ can be
chosen so that
\[
\begin{pmatrix}\alpha', & \beta'\\ \gamma', & \delta'\end{pmatrix} \equiv \begin{pmatrix}\alpha, & \beta\\ \gamma, & \delta\end{pmatrix}\ (\mathrm{mod}\ p) \ \text{and}\ \alpha'\delta'-\beta'\gamma'=1
\]
holds. All substitutions of $\mathfrak{L}_0$ are therefore contained among those considered in \S\,30.
But here we can also easily exhibit the actual
composition of an arbitrary substitution in $\mathfrak{L}_0$ out of $A$ and
$B$, and so carry out the proof independently of \S\,30.

For if in the first place $c$ is an arbitrary number not divisible by $p$,
then actual calculation easily gives
\weq{11}{A^c=\begin{pmatrix}1, & 0\\ c, & 1\end{pmatrix}, \quad C=\begin{pmatrix}c, & 0\\ 0, & c^{-1}\end{pmatrix}=A^{-c^{-1}}BA^{-c}BA^{-c^{-1}}B,}
%CONT>>
